% !TEX root = ../ApproxDA-TransUNet.tex
\begin{abstract}
Dual attention, which combines spatial and channel self-attention, is effective for medical image segmentation, but its quadratic cost has motivated windowed, low-rank, and grouped approximations that are usually treated as accuracy-neutral. How these approximations affect segmentation accuracy, and whether the effect is the same across tasks, has received little controlled study. We present ApproxDA-TransUNet, which keeps the architecture of DA-TransUNet and replaces only its attention modules with an approximation controlled by three independent parameters: spatial window size $M$, projection rank $r$, and channel group number $G$. Varying one parameter at a time on Synapse multi-organ CT, Kvasir-SEG, and ISIC~2018, we observe strongly dataset-dependent window sensitivity: the Dice range across window sizes is \tbd{2.30} points on Synapse but only \tbd{0.64} and \tbd{0.50} points on Kvasir-SEG and ISIC~2018. With an appropriately selected window, ApproxDA-TransUNet improves Dice over a DA-TransUNet re-trained under the same protocol on all three datasets, while using fewer parameters and FLOPs in its attention blocks. These results indicate that attention approximation is not accuracy-neutral and that its window scale should be treated as a task-dependent design choice.
\end{abstract}

\begin{IEEEkeywords}
medical image segmentation, dual attention, attention approximation, window size, ablation study
\end{IEEEkeywords}
